\begin{scanpage}{0098}{81}
Here $\sqrt{a/R}$ is the unfolding coefficient of the wave. It shows how the field is formed with increasing distance from the disk: the diffracted wave which is cylindrical near the disk unfolds into a spherical wave as the distance from it increases. The coefficient $(\sin\gamma+\sin\vartheta)^{-1/2}$ is proportional to the width of the luminous region on the disk or, in other words, to the width of the first Fresnel zone. Thus, in Equations (12.01) and (12.12) the functions $f^1$ and $g^1$ depend only on the body's geometry --- more precisely, on the character of the discontinuity.

Therefore, it is entirely natural to assume that the similar waves which are being scattered by a cylinder have the same structure and differ only in the functions $f^1$ and $g^1$ which correspond in this case to a rectangular wedge. Consequently, in the direction toward the source, the field from a nonuniform part of the current flowing on the cylinder may be represented when $ka\sin\vartheta\gg1$ in the following way:
\[
\left.\begin{aligned}
E_\varphi=-H_\vartheta&=\frac{ia}{\sqrt{2\pi\zeta}}E_{0x}\{f^1(1)e^{i(\zeta-3\pi/4)+ikl\cos\vartheta}\\
&\quad-[f^1(2)e^{ikl\cos\vartheta}+f^1(3)e^{-ikl\cos\vartheta}]e^{-i(\zeta-3\pi/4)}\}\frac{e^{ikR}}R,\\
E_\vartheta=H_\varphi&=\frac{ia}{\sqrt{2\pi\zeta}}H_{0x}\{g^1(1)e^{i(\zeta-3\pi/4)+ikl\cos\vartheta}\\
&\quad-[g^1(2)e^{ikl\cos\vartheta}+g^1(3)e^{-ikl\cos\vartheta}]e^{-i(\zeta-3\pi/4)}\}\frac{e^{ikR}}R.
\end{aligned}\right\}\tag{14.02}
\]
In accordance with \S\ 4, the functions $f^1$ and $g^1$ are determined by the equations
\[
\left.\begin{aligned}
\left.\begin{matrix}f^1(1)\\g^1(1)\end{matrix}\right\}
&=\frac{\sin\frac\pi n}n\left(\frac1{\cos\frac\pi n-1}\mp\frac1{\cos\frac\pi n-\cos\frac{\pi-2\vartheta}n}\right)\pm\frac{\cos\vartheta}{2\sin\vartheta},\\
\left.\begin{matrix}f^1(2)\\g^1(2)\end{matrix}\right\}
&=\frac{\sin\frac\pi n}n\left(\frac1{\cos\frac\pi n-1}\mp\frac1{\cos\frac\pi n-\cos\frac{2\vartheta}n}\right)\mp
\end{aligned}\right\}\tag{14.03}
\]
\begin{center}(Equation continued on next page.)\end{center}
\end{scanpage}
